\section{Conclusion}
\label{sec:conclusion}

We studied whether a video game can teach strangers to coordinate on unrelated tasks---the cross-domain zero-shot coordination problem.
Using LLM agents as human proxies, we designed and compared three coordination-training games across a 31-task evaluation suite spanning focal points, role assignment, and convention formation.

Our main findings are:
(1) Game training produces up to 2.9$\times$ improvement on individual focal point tasks and ${\sim}1.5\times$ mean improvement on hard tasks.
(2) The AI-designed multi-skill game \syncminds achieves the most balanced improvement, including 1.98$\times$ on hard role assignment (50.5\%$\to$100\%).
(3) Single-skill training creates a convergence--differentiation tension: focal point training improves matching tasks but destroys complementary role assignment.
Only multi-skill training resolves this tension.

These results provide proof-of-concept that AI can design coordination-enhancing games with measurable cross-domain transfer.
The key design principle is \emph{multi-skill training}: effective games must teach both when to match a partner and when to differentiate, giving players meta-cognitive awareness of which coordination strategy each situation demands.

\para{Future work.}
The most important next step is human validation.
While LLM agents enable rapid, controlled experimentation, they coordinate through shared training data rather than lived experience.
Human participant studies are needed to confirm that multi-skill game training produces real coordination gains.
Beyond validation, we see three promising directions: (1) reinforcement learning for game design optimization, (2) cross-cultural coordination studies, and (3) interactive implementation of \syncminds as a playable game for deployment at scale.
